\section{Conclusion}
\label{sec:conclusion}

We formalize gate-keeping as a capacity-constrained Bayesian signaling problem at three layers (paper, reader, researcher), parameterized by the submission volume $\Nsub$ that prior signaling models hold fixed. Analytically, fixed capacity drives mutual information and tail recall to zero in $\Nsub$ (Kendall $\tau \le -0.95$ uniformly), and fixed-capacity researcher evaluation collapses to bean counting once the field is large enough to make publication count tied at zero for almost everyone ($95.5\%$ at $M_r{=}20{,}000$). Empirically, $14$ years of \iclr data ($\Nsub: 67 \to 19{,}814$) show a log-linear decay of the standardized posterior contrast at $-0.32$ per $e$-fold of $\Nsub$ ($R^2 = 0.85$, $p = 1.1\times10^{-3}$), accompanied by a $0.86 \to 1.53$ rise in mean per-paper reviewer SD. The two patterns are jointly consistent with a {\bf load-induced noise} interpretation: \iclr expanded capacity to keep accept rate constant, but the reviewer pool did not expand proportionally, so $\sigrev$ grew and the signal degraded anyway.

\para{Key takeaway.} No capacity policy preserves all three of per-paper precision, researcher identifiability, and reader-level relevance as $\Nsub$ grows; either the signal degrades or the pool from which the signal selects becomes too large for downstream consumers to use. The venue label survives as a necessary filter and ceases to be a sufficient one.

\para{Future work.} Three directions sit naturally next to this paper. (i) An empirical ablation across venue formats (single-blind, double-blind, structured \openreview, journal-with-revisions) would let us estimate $\sigrev^2(\Nsub, \text{format})$ and test whether mechanism-design changes can break the load-induced noise channel. (ii) An adaptive capacity policy $\Kcap(\Nsub) = \alpha \log \Nsub$ trades precision and recall in a different way than the three policies we test; the analytical machinery is in hand and the sweep is one parameter. (iii) The reader-trust-collapse model deserves a closed-form treatment with venue priors and a misspecification penalty, so that the threshold at which collapse becomes generic can be characterized in terms of $\Nsub$, $\sigrev$, and reader expectations rather than a numerical $\tau_{\text{bad}}$.
